% ===== BOT 08: final_prune_structgs — prune cuối theo pruning score & multi-view consistency =====

% --- Frame 1: Vị trí trong pipeline ---
\begin{frame}{\texttt{final\_prune\_structgs}: tỉa cuối cùng nằm ở đâu?}
\scriptsize
\begin{itemize}
    \item SADGS có \textbf{hai tầng tỉa} tách biệt: (1) tỉa \emph{trong} vòng mật độ hoá, gọi mỗi \texttt{densification\_interval} từ \texttt{densify\_and\_prune\_structgs} (\texttt{gaussian\_model.py:1012--1046}); (2) tỉa \emph{cuối giai đoạn} bằng \texttt{final\_prune\_structgs} (\texttt{gaussian\_model.py:1059--1066}).
    \item Tầng (2) được thiết kế để chạy tại các mốc \texttt{prune\_iterations}, mặc định \texttt{[4000, 8000]} (\texttt{train.py:529}), được truyền thẳng vào \texttt{training(...)} như tham số cuối (\texttt{train.py:38}, \texttt{train.py:555}).
    \item Mốc kiểm tra: \texttt{if iteration in prune\_iterations:} (\texttt{train.py:412}). Tại đó code lấy toàn bộ camera huấn luyện rồi lọc qua \texttt{sampling\_cameras} (\texttt{train.py:413--414}, \texttt{utils/freq\_utils.py:12}, mặc định \texttt{mode="fps", num\_cams=60}) -- tức đã chuẩn bị sẵn tập view để đánh giá nhất quán đa góc nhìn.
    \item \alert{Nhưng} hai dòng dùng tập view đó lại đang bị comment (\texttt{train.py:418--419}) -- xem Frame 5.
\end{itemize}
\centering
\includegraphics[width=0.62\linewidth]{sadgsx/08_timeline.png}
\captionof{figure}{\footnotesize Dòng thời gian huấn luyện 30k iteration: vùng mật độ hoá 500--15\,000, hai mốc \texttt{prune\_iterations} 4000 và 8000 (\texttt{train.py:412, 529}; \texttt{arguments/\_\_init\_\_.py:75, 91--92}).}
\end{frame}

% --- Frame 2: Đọc code từng dòng ---
\begin{frame}[fragile]{Đọc từng dòng: công thức tỉa là một phép OR hai điều kiện}
\scriptsize
\begin{block}{\texttt{scene/gaussian\_model.py:1059--1066}}
\begin{semiverbatim}
def final_prune_structgs(self, min_opacity, pruning_score = None):   # :1059
    prune_mask  = (self.get_opacity < min_opacity).squeeze()         # :1063
    scores_mask = pruning_score > 0.9                                # :1064
    final_prune = torch.logical_or(prune_mask, scores_mask)          # :1065
    self.prune_points(final_prune)                                   # :1066
\end{semiverbatim}
\end{block}
\footnotesize
Toàn bộ hàm chỉ gồm một công thức mặt nạ, với $\alpha_i$ là opacity sau sigmoid và $s_i$ là \texttt{pruning\_score}:
\[
\text{prune}_i \;=\; \underbrace{\big[\alpha_i < \tau_\alpha\big]}_{\text{opacity yếu}} \;\vee\; \underbrace{\big[s_i > \tau_s\big]}_{\text{bất nhất đa view}},\qquad \tau_\alpha = \texttt{min\_opacity},\quad \tau_s = 0{,}9 \text{ (hằng số cứng)}
\]
\begin{itemize}\scriptsize
    \item $\tau_s = 0{,}9$ \textbf{hard-code} ngay trong thân hàm (\texttt{:1064}) -- không có tham số dòng lệnh nào chỉnh được, khác hẳn \texttt{prune\_ratio\_threshold} $=0{,}8$ vốn khai báo trong \texttt{arguments/\_\_init\_\_.py:149}.
    \item $\tau_\alpha$ do lời gọi quyết định; lời gọi duy nhất trong repo dùng \texttt{min\_opacity = 0.1} (\texttt{train.py:419}).
    \item Mặt nạ đi thẳng vào \texttt{prune\_points} (\texttt{gaussian\_model.py:528}), hàm này cắt đồng bộ \emph{cả} tensor tham số lẫn mọi bộ tích luỹ: \texttt{xyz\_gradient\_accum}, \texttt{max\_radii2D}, \texttt{filter\_3D}, \texttt{accum\_eta}, \texttt{accum\_view\_count}, \texttt{eta\_high/mid/low\_count} (\texttt{:540--560}).
\end{itemize}
\end{frame}

% --- Frame 3: Ngưỡng cắt, minh hoạ phân bố ---
\begin{frame}{Ngưỡng cắt $\tau_s=0{,}9$ và $\tau_\alpha=0{,}1$ trên phân bố thực tế}
\footnotesize
\begin{columns}[T]
\begin{column}{0.56\textwidth}
\centering
\includegraphics[width=\linewidth]{sadgsx/08_score_hist.png}
\end{column}
\begin{column}{0.42\textwidth}
\scriptsize
\begin{itemize}
    \item Vế trái: $[s_i > 0{,}9]$ chỉ bắt \emph{đuôi phải} của phân bố -- đây là thiết kế \textbf{bảo thủ}: chỉ Gaussian bất nhất ở gần như mọi view mới bị loại.
    \item Vế phải: $[\alpha_i < 0{,}1]$ là ngưỡng \textbf{rất mạnh} so với 3DGS gốc. Code ghi rõ trong comment \texttt{train.py:364}: \texttt{min\_opacity=0.1, \# 0.005 is a good default} -- tức SADGS chủ động dùng ngưỡng cao gấp 20 lần mặc định.
    \item Vì là phép \textbf{OR}, tổng số Gaussian bị loại $\ge$ số bị loại bởi từng vế. Trên phân bố minh hoạ ($N=20\,000$): $\approx 5{,}6\%$ bị loại, trong đó nhánh opacity $\approx 1{,}0\%$ và nhánh nhất quán $\approx 4{,}7\%$.
    \item Ước lượng: ở quy mô $\sim 1$--$3$ triệu Gaussian điển hình, mỗi mốc prune loại bậc $10^4$--$10^5$ điểm.
\end{itemize}
\end{column}
\end{columns}
\captionof{figure}{\footnotesize Hai nhánh của $\text{prune}_i=[\alpha_i<0{,}1]\vee[s_i>0{,}9]$ (\texttt{gaussian\_model.py:1063--1066}).}
\end{frame}

% --- Frame 4: pruning_score đến từ đâu ---
\begin{frame}{\texttt{pruning\_score} được truyền từ đâu? -- và điều gì xảy ra nếu \texttt{None}}
\footnotesize
\begin{columns}[T]
\begin{column}{0.52\textwidth}
\scriptsize
\begin{itemize}
    \item Nguồn dự kiến duy nhất là dòng đã bị comment: \\ \texttt{\_, pruning\_score, \_, \_, \_, \_ = compute\_gaussian\_score\_structgs(camlist, gaussians, pipe, bg, opt)} (\texttt{train.py:418}) -- một hàm trả về 6 giá trị, tính trên \texttt{camlist} gồm 60 camera lấy bằng FPS.
    \item \alert{Phát hiện quan trọng:} \texttt{grep -rn "compute\_gaussian\_score\_structgs"} trên toàn bộ repo chỉ trả về \textbf{đúng 1 kết quả} -- chính dòng comment đó. Hàm \textbf{không được định nghĩa và không được import} ở bất kỳ đâu. Bỏ comment ngay sẽ lỗi \texttt{NameError}.
    \item Chữ ký hàm khai báo \texttt{pruning\_score = None} (\texttt{:1059}), nhưng thân hàm \textbf{không kiểm tra None}: dòng \texttt{:1064} thực hiện \texttt{None > 0.9} $\Rightarrow$ \texttt{TypeError}. Tức giá trị mặc định là một cái bẫy, hàm bắt buộc phải nhận tensor thật.
    \item Tín hiệu nhất quán đa view \emph{đang} được đo ở nơi khác: \texttt{low\_ratio}$_i=$\texttt{eta\_low\_count}$_i/$\texttt{accum\_view\_count}$_i$ (\texttt{train.py:341--342}), dùng cho prune trong vòng densify.
\end{itemize}
\end{column}
\begin{column}{0.46\textwidth}
\centering
\includegraphics[width=\linewidth]{sadgsx/08_mvconsistency.png}
\end{column}
\end{columns}
\captionof{figure}{\footnotesize Ý nghĩa $s_i$: gộp bất nhất tái dựng trên 60 view (\texttt{sampling\_cameras}, \texttt{freq\_utils.py:12}) thành 1 điểm số; chỉ hàng vượt $\tau_s=0{,}9$ mới bị cắt.}
\end{frame}

% --- Frame 5: So sánh với prune trong vòng densify ---
\begin{frame}{Khác biệt với phép prune trong vòng mật độ hoá}
\scriptsize
\begin{table}
\centering
\begin{tabular}{p{2.4cm}p{4.4cm}p{4.4cm}}
\toprule
 & \textbf{Prune trong vòng densify} \newline \texttt{densify\_and\_prune\_structgs} & \textbf{Prune cuối} \newline \texttt{final\_prune\_structgs} \\
\midrule
Vị trí code & \texttt{gaussian\_model.py:1012--1046} & \texttt{gaussian\_model.py:1059--1066} \\
Tần suất & mỗi \texttt{densification\_interval}, trong 500--15\,000 & chỉ tại \texttt{prune\_iterations} $=[4000,8000]$ \\
Điều kiện & $[\alpha<0{,}1]\vee[r_{2D}>\texttt{max\_screen}]\vee[\max s>0{,}1\cdot\text{extent}]\vee$ \texttt{custom\_prune\_mask} & $[\alpha<0{,}1]\vee[s>0{,}9]$ -- \emph{chỉ} opacity và điểm số đa view \\
Tín hiệu đa view & \texttt{low\_ratio}$>$\texttt{prune\_ratio\_threshold}$=0{,}8$ (\texttt{train.py:354}, \texttt{arguments:149}), reset bộ đếm sau mỗi vòng & \texttt{pruning\_score} tính lại \emph{một lần} trên 60 view FPS, không phụ thuộc bộ đếm đang tích luỹ \\
Khi có score & \alert{ngẫu nhiên hoá}: \texttt{remove\_budget = int(0.5*to\_remove)}, lấy mẫu \texttt{torch.multinomial} theo $1/(10^{-6}+1-s)$ rồi AND với \texttt{prune\_mask} (\texttt{:1029--1042}) $\Rightarrow$ chỉ bỏ \textbf{một nửa} ứng viên & \textbf{tất định}: OR rồi cắt hết, không có budget, không lấy mẫu \\
Hậu xử lý & ép $\alpha \leftarrow \min(\alpha, 0{,}8)$ qua \texttt{inverse\_sigmoid} (\texttt{:1048--1050}) & không có -- chỉ \texttt{prune\_points} \\
\bottomrule
\end{tabular}
\end{table}
\footnotesize
$\Rightarrow$ Tầng trong vòng densify là \emph{tỉa tăng dần, có phanh} (giữ lại 50\% ứng viên); tầng cuối là \emph{tỉa dứt khoát} nhằm dọn floater còn sót.
\end{frame}

% --- Frame 6: Trạng thái comment out ---
\begin{frame}[fragile]{Trạng thái thực tế: lời gọi \textbf{đang bị comment}}
\scriptsize
\begin{block}{\texttt{train.py:412--419} -- nguyên văn}
\begin{semiverbatim}
if iteration in prune_iterations:                                          # :412
    my_viewpoint_stack = scene.getTrainCameras().copy()                    # :413
    camlist = sampling_cameras(my_viewpoint_stack)                         # :414
    prune_mask = (gaussians.get_opacity < 0.1).squeeze()                   # :416
    gaussians.prune_points(prune_mask)                                     # :417
    # _, pruning_score, ... = compute_gaussian_score_structgs(...)         # :418
    # gaussians.final_prune_structgs(min_opacity=0.1, pruning_score=...)   # :419
\end{semiverbatim}
\end{block}
\footnotesize
\begin{itemize}\scriptsize
    \item Cái \emph{đang chạy} tại iteration 4000 và 8000 chỉ là nhánh opacity thuần: $\text{prune}_i=[\alpha_i<0{,}1]$. Toàn bộ vế $[s_i>0{,}9]$ biến mất.
    \item \texttt{camlist} ở dòng \texttt{:414} vẫn được tính nhưng \textbf{không ai dùng} -- dấu vết rõ ràng rằng nhánh multi-view đã bị tắt chứ không phải chưa từng có.
    \item Đây không phải trường hợp cá biệt: nhánh \texttt{pruning\_score} trong vòng densify cũng bị vô hiệu vì lời gọi truyền \texttt{pruning\_score=None} (\texttt{train.py:369}) $\Rightarrow$ khối \texttt{multinomial} \texttt{:1029--1042} không bao giờ chạy, code rơi vào \texttt{else: self.prune\_points(prune\_mask)} (\texttt{:1044}).
    \item Hệ quả: \textbf{cả hai} đường dùng \texttt{pruning\_score} trong repo đều đang chết. Tín hiệu nhất quán đa view duy nhất còn sống là \texttt{low\_ratio} qua \texttt{custom\_prune\_mask}.
\end{itemize}
\end{frame}

% --- Frame 7: Hệ quả về kích thước mô hình ---
\begin{frame}{Hệ quả: mô hình cuối to hơn mức cần thiết}
\footnotesize
\begin{columns}[T]
\begin{column}{0.54\textwidth}
\centering
\includegraphics[width=\linewidth]{sadgsx/08_modelsize.png}
\end{column}
\begin{column}{0.44\textwidth}
\scriptsize
\begin{itemize}
    \item Mỗi Gaussian lưu 62 float32 trong \texttt{.ply} kiểu 3DGS (xyz, normal, 3 hệ số SH bậc 0, 45 hệ số SH bậc cao, opacity, 3 scale, 4 quaternion) $\Rightarrow$ $\approx 248$ byte/điểm.
    \item Bỏ mất vế $[s_i>0{,}9]$ nghĩa là các \textbf{floater} -- Gaussian có opacity cao nhưng chỉ đúng ở một vài góc nhìn -- \emph{không bị vế nào bắt}: chúng vượt qua $\tau_\alpha=0{,}1$ một cách thoải mái.
    \item Floater gây hại kép: (i) phình dung lượng \texttt{.ply} và giảm FPS render; (ii) tạo vệt sương/đốm khi camera rời khỏi quỹ đạo train -- đúng lỗi mà tỉa theo nhất quán đa view sinh ra để khử.
    \item Với mô hình digital twin cần xem tự do quanh vật thể, đây là mất mát chất lượng \emph{chỉ thấy ở view mới}, không thấy trên PSNR tập train.
\end{itemize}
\end{column}
\end{columns}
\captionof{figure}{\footnotesize So sánh kích thước mô hình: không tỉa cuối / chỉ $\alpha<0{,}1$ (code hiện tại) / đủ $[\alpha<0{,}1]\vee[s>0{,}9]$.}
\end{frame}

% --- Frame 8: Đánh đổi ngưỡng ---
\begin{frame}{Đánh đổi: hạ $\tau_s$ tỉa mạnh hơn -- đổi lấy gì?}
\footnotesize
\begin{columns}[T]
\begin{column}{0.52\textwidth}
\centering
\includegraphics[width=\linewidth]{sadgsx/08_tradeoff.png}
\end{column}
\begin{column}{0.46\textwidth}
\scriptsize
Quét $\tau_s$ trong công thức $\text{prune}_i=[\alpha_i<\tau_\alpha]\vee[s_i>\tau_s]$:
\begin{itemize}
    \item $\tau_s\to 1$: gần như không tỉa theo đa view -- tương đương trạng thái hiện tại của code.
    \item $\tau_s=0{,}9$ (giá trị hard-code): giữ $\approx 94\%$ Gaussian, cắt gọn đuôi bất nhất mà gần như không đụng phần thân phân bố $\Rightarrow$ chất lượng bão hoà, dung lượng giảm.
    \item $\tau_s<0{,}7$: bắt đầu ăn vào vùng Gaussian hợp lệ; số điểm giữ lại rơi nhanh và chất lượng tụt -- vùng tỉa quá tay.
    \item Đường PSNR trong hình là \textbf{mô phỏng dạng bão hoà} để minh hoạ hình dạng đánh đổi, không phải số đo từ repo (repo không chứa log benchmark cho nhánh này).
\end{itemize}
\end{column}
\end{columns}
\captionof{figure}{\footnotesize Tỉ lệ Gaussian giữ lại (trục trái) và chất lượng mô phỏng (trục phải) theo $\tau_s$; vạch chấm là $\tau_s=0{,}9$ trong \texttt{gaussian\_model.py:1064}.}
\end{frame}

% --- Frame 9: Tổng kết & khuyến nghị ---
\begin{frame}{Tổng kết \texttt{final\_prune\_structgs}}
\footnotesize
\begin{enumerate}
    \item \textbf{Công thức:} $\text{prune}_i=[\alpha_i<\texttt{min\_opacity}]\vee[s_i>0{,}9]$ -- một phép OR tất định, không budget, không lấy mẫu ngẫu nhiên (\texttt{gaussian\_model.py:1063--1066}).
    \item \textbf{Tham số:} $\tau_\alpha$ truyền vào ($=0{,}1$ tại lời gọi), $\tau_s=0{,}9$ hard-code; mốc chạy \texttt{prune\_iterations} mặc định \texttt{[4000, 8000]} (\texttt{train.py:529}).
    \item \textbf{Trạng thái:} lời gọi bị comment (\texttt{train.py:419}); nguồn điểm số \texttt{compute\_gaussian\_score\_structgs} \alert{không tồn tại trong repo}; \texttt{pruning\_score=None} mặc định sẽ gây \texttt{TypeError} vì hàm không kiểm tra None.
    \item \textbf{Đang chạy thay thế:} \texttt{prune\_points}$([\alpha<0{,}1])$ tại cùng hai mốc (\texttt{train.py:416--417}) -- giữ nhánh opacity, mất nhánh nhất quán đa view.
    \item \textbf{Để khôi phục an toàn:} (a) cài đặt \texttt{compute\_gaussian\_score\_structgs} trả về $s\in[0,1]$ theo \emph{độ bất nhất} (giá trị lớn = xấu, khớp hướng so sánh \texttt{>} ở \texttt{:1064}); (b) thêm \texttt{if pruning\_score is None: scores\_mask = torch.zeros\_like(prune\_mask)}; (c) đưa $0{,}9$ ra \texttt{arguments/\_\_init\_\_.py} cạnh \texttt{prune\_ratio\_threshold}.
\end{enumerate}
\end{frame}
